\paragraph{Missing step.} The full target of Problem~C is the gauge-fixed quadratic scalar operator on the locked vacuum in a constructed frame: the canonically normalised kinetic coefficient of the pure-trace scalar, including the induced Einstein--Hilbert term, its mixing with \(T_{\mu\nu}\), its non-minimal curvature coupling, its matter coupling, and its pole mass as a derived output, with Assumption~\ref{ass:gap-psi} discharged. On the bare data of Presentation~A this operator is not yet definable: three of the inputs it needs are among the things the target is meant to establish. (i)~A frame. The transformation law of Equation~\eqref{eq:Phi-transform} is a field-content choice of Presentation~B and is not derived from Presentation~A alone; the stated Repair~I map to an orthonormal frame is obstructed above; the remark following Definition~\ref{def:vierbein} records that a corrected Repair~I would need additional background tensors or an independent coframe, which is part of Problem~C; and the route actually used posits the frame through Assumptions~\ref{ass:nondeg} and~\ref{ass:frame-lock}. (ii)~\(G_{\mathrm{eff}}\). The coefficient \(Z\) of Equation~\eqref{eq:Z-sigma} is written in terms of \(G_{\mathrm{eff}}\), whose derivation is Problem~F, so the scalar wave-function renormalisation ties Problem~C to Problem~F. (iii)~The constraint structure of the massless sector. The trace-constrained Fierz--Pauli structure of the quadratic tensor block is not shown to follow from the NAMCG quadratic block, and whether the \(\psi\) sector of Chapter~\ref{ch:massless-tensor} supplies it is unresolved and is a dependency of Problem~C; the \(\psi\) of Chapter~\ref{ch:massless-tensor} is the trace of the tensor fluctuation, distinct from the mode labels of \(K\) used in the Problem~F entry. With these three inputs supplied, the recorded results of this section (the tree radial mass of Lemma~\ref{lem:msigma}, the wave-function factor \(Z\) of Equation~\eqref{eq:Z-sigma} at the stated truncation, and the off-shell trace--tensor mixing term) are computations within posited frame data; the loop-level kinetic coefficient, the mixing, the non-minimal coupling, the matter coupling and the pole mass are not outputs of the matrix theory, and in the conjugate-coordinate class the \(q=0\) non-derivative sector of the trace channel carries a bounded obstruction recorded in the Problem~F entry.
%---SPLIT---%
\emph{Trace-channel tadpole and the non-derivative sector at \(q=0\).} Under the hypotheses of the preceding set of paragraphs (equal weights, \(V=T\), frozen symbol, product background, static channel \(q\cdot v=0\), one loop), the \(q=0\) part of every density term that enters the trace channel carrying \(\sqrt{|g|}\) alone is fixed by the determinant. Write such a term as \(W_i\sqrt{|g|}\) at the background and take the channel along the pure-trace direction \(h_{\mu\nu}=(t/4)\bar g_{\mu\nu}\), \(t=h^\mu{}_\mu\): then \(\sqrt{|g+h|}=\sqrt{|g|}\,(1+t/2+t^2/16+O(t^3))\), by the expansion of \(\sqrt{\det(1+\epsilon)}\) to second order, so the whole \(q=0\) contribution of the term to the reduced \(t\)-form runs along one profile common to all such terms. Its tadpole \(T_i\) (the coefficient of \(t\) at \(q=0\)) and its contribution \(\mathcal M_i\) to the term \(\tfrac12\mathcal Mt^2\) are then proportional with a fixed ratio, invariant under the choice of continuation, while their common sign follows it: with the formal continuation \(S_E=-S_L\), \(T_i=-\sqrt{|g|}W_i/2\) and \(\mathcal M_i=-\sqrt{|g|}W_i/8\), and with \(q_*^2=2\mathcal ML\) as above, \(q_*^2=\tfrac12L\sum_iT_i\). Two consequences. First, no supplied set of terms carrying \(\sqrt{|g|}\) alone, including a cosmological-type counterterm, cancels the \(q=0\) tadpole and leaves a nonzero \(\mathcal M\): the cancellation condition \(\sum_iW_i=0\) cancels the quadratic part with it, so with the recorded terms a background flat at \(q=0\) has \(q_*^2=0\) and the pole of the vacuum-term display sits at \(q^2=q_b^2\), which for \(\hat U>0\) has the tachyonic sign recorded above. Conversely the Yukawa threshold \(q_*^2=q_b^2\) at \(\hat U>0\) requires an uncancelled tadpole of size \(|T|=F\hat U\chi^4/(32\pi^2m^8)\) in the same convention. Second, taken alone and not cancelled, \(U_4\) contributes \(\mathcal M=F\sqrt{|g|}\hat U\chi^4/256\pi^2\), of the sign that pulls \(q_*^2\) positive for \(\hat U>0\), and \(q_*^2=6\sqrt{|g|}\hat U\chi^2m^6/(\hat V+2\hat T)\); its ratio to \(q_b^2\) is \(\sqrt{|g|}m^8/2\), and the comparison of magnitudes turns on the conventions for \(m\) and \(\sqrt{|g|}\) and is not settled here. The sign of the quadratic part of the determinant profile is direction-dependent (along a rank-one direction the second-order part is \(-\lambda^2/8\)), so the pure-trace channel used by the reduced form is part of the statement. The derivative in \(\chi\) that sets \(\beta_K\) survives these cancellations, because the tree density carries \(\rho_M\) and not \(\sqrt{|g|}\) as recorded (and no identification of \(\rho_M\) with a determinant of the fluctuating frame is made in this class), so it supplies no \(t\)-channel mass: the requirement recorded above, that the vacuum term cancel at the minimum without cancelling its derivative in \(\chi\), is compatible with these identities, and the surviving \(\chi\) derivative is not the missing \(\mathcal M\). Within the recorded class, then, the \(q=0\) terms flagged above as not computed cannot be cancelled by adding further vacuum-energy terms carrying \(\sqrt{|g|}\) alone; the candidate that could supply the required \(\mathcal M\) without cancelling its own tadpole is the response of the fibre levels to \(\bar h\) in \(v^M-2\sigma^M\), which is recorded as not computed. These statements are for the hypotheses of the preceding set of paragraphs and the channel convention stated here; they do not compute the fibre-level response, the loop kinetic term of \(\chi\), the second-order non-derivative terms at \(\chi^2\lesssim\vartheta\), the rotating class or higher loops, do not settle the convention-bound magnitudes, do not identify \(t\) with a quantity of the emergent metric, and give no Cassini or Newtonian comparison. Problems~C, F and~N remain open and PROGRAMME-CONDITIONAL.
